\documentclass[10pt,a4paper]{extarticle}

% ── Packages ──────────────────────────────────────────────────────────
\usepackage[utf8]{inputenc}
\usepackage[T1]{fontenc}
\usepackage{lmodern}
\usepackage[top=1.8cm, bottom=1.8cm, left=1.7cm, right=1.7cm]{geometry}
\usepackage{graphicx}
\usepackage{booktabs}
\usepackage{array}
\usepackage{tabularx}
\usepackage{multirow}
\usepackage{longtable}
\usepackage{amsmath}
\usepackage[round,authoryear]{natbib}
\usepackage{hyperref}
\usepackage{xcolor}
\usepackage{enumitem}
\usepackage{caption}
\usepackage{float}
\usepackage{parskip}
\usepackage{setspace}
\usepackage{fancyhdr}
\usepackage{titlesec}
\usepackage{pdflscape}
\usepackage[most]{tcolorbox}

% ── Colours ───────────────────────────────────────────────────────────
\definecolor{imperialblue}{RGB}{0,62,116}

% ── Hyperlinks ────────────────────────────────────────────────────────
\hypersetup{
    colorlinks=true,
    linkcolor=imperialblue,
    urlcolor=imperialblue,
    citecolor=imperialblue
}

% ── Spacing ───────────────────────────────────────────────────────────
\setstretch{1.15}

% ── Section formatting ────────────────────────────────────────────────
\titleformat{\section}{\Large\bfseries}{}{0em}{}
\titleformat{\subsection}{\large\bfseries}{\thesubsection}{0.5em}{}
\titlespacing*{\section}{0pt}{1.5em}{0.8em}
\titleformat{\subsubsection}{\normalsize\bfseries}{\thesubsubsection}{0.5em}{}
\titlespacing*{\subsection}{0pt}{1.2em}{0.6em}
\titlespacing*{\subsubsection}{0pt}{1.0em}{0.4em}

% ── Header / Footer ──────────────────────────────────────────────────
\pagestyle{fancy}
\fancyhf{}
\fancyhead[L]{}
\fancyhead[R]{}
\fancyfoot[R]{\thepage}
\renewcommand{\headrulewidth}{0pt}

% ══════════════════════════════════════════════════════════════════════
\raggedbottom
\begin{document}

% ── Title Page ────────────────────────────────────────────────────────
\thispagestyle{empty}
\begin{center}
\vspace*{2cm}

{\LARGE\bfseries ELEC70126: Internet of Things and Applications}

\vspace{0.5cm}
{\large Imperial College London}

\vspace{2cm}

\newsavebox{\infotablebox}
\savebox{\infotablebox}{%
\begin{tabular}{@{} >{\bfseries}l l @{}}
\toprule
Student CID & \textit{06053600} \\
\midrule
Student Name & \textit{Yasmin Akhmedova} \\
\midrule
Assessment Due & \textit{2pm 25/03/2026} \\
\midrule
Study Topic & \textit{Exploration of Sleeping Environment Conditions} \\
            & \textit{for Improved Sleep Quality} \\
\midrule
GitHub Project Repository (Detailed Coding Report) & \textit{\href{https://github.com/ya325icacuk/iot-sleep-environment-coursework}{ya325icacuk/iot-sleep-environment-coursework}} \\
\midrule
Interactive Visualisation Dashboard & \textit{\href{https://sleep-monitor.streamlit.app}{sleep-monitor.streamlit.app}} \\
\bottomrule
\end{tabular}}
\usebox{\infotablebox}

\vspace{1.5cm}

{\large\bfseries Contents}

\vspace{0.4cm}

\newlength{\tocwidth}
\setlength{\tocwidth}{\wd\infotablebox}
\noindent\begin{minipage}{\tocwidth}
\makeatletter
\newcommand{\tocline}[2]{#1\nobreak\leaders\hbox to 0.5em{\hss.\hss}\hfill\nobreak\hbox to 1.5em{\hfil#2}\par\vspace{4pt}}
\makeatother
\tocline{Executive Summary}{2}
\tocline{Part 1: Sensing}{2}
\tocline{Part 2: Internet of Things}{6}
\tocline{Conclusion}{11}
\tocline{References}{12}
\tocline{Appendix: Figures and Illustrations}{13}
\end{minipage}

\vspace{\fill}
\end{center}
\newpage

% ── Executive Summary ─────────────────────────────────────────────────
\section*{Executive Summary}
Sleep quality is a growing public health concern, with a global survey finding that people report losing nearly three nights of restorative sleep each week~\citep{resmed2025}. Multiple environmental factors are known to affect sleep~\citep{caddick2018,yasmeen2025}. However, few studies combine low-cost sensor data with consumer wearable sleep tracking in real home settings~\citep{basner2023}. This report addresses that gap by investigating which bedroom environment conditions are most strongly associated with poor sleep quality, using affordable IoT hardware costing under \pounds30 in a central London bedroom.

Over 14 consecutive nights (9 to 22 February 2026), three data sources were collected: a custom sensor system logging temperature, humidity, light, and sound, hourly air quality data from the Breathe London API, and daily sleep metrics from an Ultrahuman Ring wearable. All data was filtered to a consistent 11\,pm to 9\,am sleep window and analysed using four methods: Spearman correlation to identify which factors relate to sleep quality, OLS regression to test their combined effects, a median-split comparison of the seven best versus seven worst nights, and minute-by-minute comparison of individual nights to examine how conditions changed across the night.

Across all four analyses, noise variability (the standard deviation of sound readings) had the strongest statistical association with sleep quality, with a Spearman correlation of $r = -0.80$ ($p < 0.01$). A regression model combining noise variability with humidity explained 64.3\% of the variance in sleep score, though noise variability was the only predictor that reached statistical significance in that model. Humidity, mean noise level, and temperature range also showed statistically significant Spearman correlations with sleep score individually. Light, PM2.5, and NO\textsubscript{2} showed no significant associations. The median-split and intra-night analyses confirmed these patterns, showing that poor sleep nights were consistently noisier, more humid, and characterised by repeated noise spikes rather than sustained loudness.

\begin{tcolorbox}[
  colback=imperialblue!5,
  colframe=imperialblue,
  leftrule=3pt,
  rightrule=0pt,
  toprule=0pt,
  bottomrule=0pt,
  boxsep=4pt,
  left=8pt,
  right=8pt,
  sharp corners
]
Based on the above findings, reducing \textbf{noise variability} through strategies such as \textbf{earplugs, white noise, or improved window insulation} is the most practical change suggested by this dataset. The findings are based on one participant over 14 nights and should be treated as indicative rather than generalisable.
\end{tcolorbox}

% ══════════════════════════════════════════════════════════════════════
\section*{Part 1: Sensing}
\setcounter{section}{1}
\setcounter{subsection}{0}
% ══════════════════════════════════════════════════════════════════════

\subsection{Introduction and Objectives}

Sleep quality has a well-documented impact on cognitive performance, physical health, and daily productivity~\citep{medic2017}. A range of bedroom environment factors can disrupt sleep, and they do so through different mechanisms. Noise, particularly from road traffic, affects sleep architecture and triggers physiological stress responses~\citep{muzet2007,halperin2014}. Temperature and humidity influence sleep through thermoregulation, since core body cooling is a critical step in falling asleep and high bedroom temperature or humidity can interfere with this process~\citep{caddick2018}. Light exposure suppresses melatonin secretion and shifts circadian timing~\citep{yasmeen2025}. Outdoor air pollutants including PM2.5 and NO\textsubscript{2} have also been associated with increased risk of sleep disorders~\citep{zhao2024}. While each of these links is individually documented, studies that objectively assess multiple bedroom environment variables alongside sleep in real home settings remain scarce~\citep{basner2023}, and fewer still combine low-cost sensor data with consumer wearable sleep tracking.

\begin{tcolorbox}[
  colback=imperialblue!5,
  colframe=imperialblue,
  leftrule=3pt,
  rightrule=0pt,
  toprule=0pt,
  bottomrule=0pt,
  boxsep=4pt,
  left=8pt,
  right=8pt,
  sharp corners
]
This study investigates the relationship between bedroom environment conditions and sleep quality over 2 weeks and has 3 objectives:
\begin{enumerate}[nosep]
    \item To \textbf{identify} which bedroom environment conditions are most strongly associated with sleep quality.
    \item To \textbf{explore} whether these associations hold when multiple conditions are considered together.
    \item To \textbf{compare} how bedroom conditions differ between good and poor sleep nights.
\end{enumerate}
\end{tcolorbox}

Two factors make this study feasible. First, low-cost IoT hardware now makes personal environment monitoring even more feasible without specialist equipment. A microcontroller board and a handful of sensors costing under \pounds30 can continuously log temperature, humidity, light, and noise throughout the night. Second, the study location, a central London flat facing a busy road, provides a naturalistic setting where noise, air quality, and other conditions vary meaningfully from night to night. Together, these make it possible to monitor multiple bedroom environment variables alongside sleep quality in a real home setting, an area where relatively few studies have been conducted. \textit{While the study is based on a single participant over 14 nights, it provides suggestive results on which bedroom conditions matter most for sleep and demonstrates an affordable, replicable approach to home-based sleep environment monitoring.}

% ──────────────────────────────────────────────────────────────────────
\subsection{Data Sources and Sensing Set-up}

To address the study's objectives, three time-series data sources of different nature were integrated. All data was collected or filtered to a consistent sleep window of 11\,pm to 9\,am over 14 consecutive nights (9 to 22 February 2026). Table~\ref{tab:sources} summarises each source, its collection method, and measured parameters.

\begin{table}[H]
\centering
\caption{Data Sources Summary}
\label{tab:sources}
\scriptsize
\begin{tabularx}{\textwidth}{c l X X c}
\toprule
\textbf{\#} & \textbf{Data Source} & \textbf{Collection Method} & \textbf{Parameters} & \textbf{Frequency} \\
\midrule
1 & Bedroom sensors & Heltec WiFi LoRa 32 V3 with DHT11, photoresistor, KY-038 mic via USB serial & Temperature, humidity, light, sound average, sound peak & Every minute \\
\midrule
2 & Air quality & Horseferry Road sensor (BL0046) via Breathe London REST API & NO\textsubscript{2} (\textmu g/m\textsuperscript{3}), PM2.5 (\textmu g/m\textsuperscript{3}) & Hourly \\
\midrule
3 & Sleep metrics & Ultrahuman Ring Air wearable via manual CSV export & Sleep score, total sleep, awake time, deep sleep, REM sleep, light sleep & Daily \\
\bottomrule
\end{tabularx}
\end{table}

\subsubsection{Sensor Hardware and Signal Processing}

\textbf{Source 1: Bedroom environment sensors.} The Heltec WiFi LoRa 32 V3 development board (ESP32-S3 microcontroller) was chosen as it provides analog-to-digital conversion pins required by the KY-038 sound sensor, as well as digital GPIO for the other sensors. Three sensors were connected, all powered at 3.3\,V to remain within the safe operating range of the ESP32-S3. Table~\ref{tab:sensorspecs} details each sensor's specifications and signal processing.

\begin{table}[H]
\centering
\caption{Sensor Specifications and Signal Processing}
\label{tab:sensorspecs}
\scriptsize
\begin{tabularx}{\textwidth}{l l X}
\toprule
\textbf{Sensor} & \textbf{Function} & \textbf{Key Specifications} \\
\midrule
KY-015 / DHT11 & Temperature, humidity & \textbullet{} $\pm2$\,\textdegree C, $\pm5$\% relative humidity accuracy \newline \textbullet{} Reads once per minute \newline \textbullet{} Retries up to 3 times with 500\,ms delays on failed reads \\
\midrule
ELB0604 photoresistor & Light detection & \textbullet{} Non-linear analog response converted to binary output (light on/off) by the onboard potentiometer, which acts as a manually adjustable threshold comparator \newline \textbullet{} Potentiometer manually calibrated to bedroom ambient light conditions \newline \textbullet{} Module outputs 0 for light and 1 for dark; inverted in the Arduino sketch so that 1 = light detected, which is more intuitive \\
\midrule
KY-038 microphone & Sound level & \textbullet{} Outputs a voltage proportional to overall sound level, with no frequency or decibel calibration \newline \textbullet{} Raw analog readings on a 0--4095 scale \newline \textbullet{} Sampled every 5\,ms ($\sim$12,000 samples per minute) \\
\bottomrule
\end{tabularx}
\end{table}

The Arduino sketch computes both sound average and sound peak each minute, and both are logged to the CSV. During analysis, only sound average was retained as it represents the sustained noise level during each one-minute interval, which is more relevant to sleep disruption than the momentary spikes captured by sound peak.

\textbf{Source 2: External air quality.} The Horseferry Road sensor (site code BL0046) was selected from the Breathe London network due to its proximity to the study location.

\textbf{Source 3: Personal sleep metrics.} Sleep data was collected using an Ultrahuman Ring Air wearable device worn each night. The ring uses photoplethysmography (PPG) and accelerometer sensors to track sleep stages and produces a daily summary including sleep score (0--100), total sleep, awake time, deep sleep, REM sleep, and light sleep (all in minutes). The sleep score is a proprietary composite metric computed by Ultrahuman's algorithm. The exact formula is not published, so it is treated as a ``black box'' measure in this study. To aid interpretation in the dashboard, sleep scores are grouped into three quality tiers: green ($\geq$80), gold (70--79), and red ($<$70).

\subsubsection{Physical Installation and Sensor Limitations}

\textbf{Physical installation.} The Heltec board and all three sensors were mounted on a breadboard placed on the bedside table, approximately 0.5\,m from the pillow and 2\,m from the bedroom window facing Horse Ferry Road. The microphone was oriented towards the window (the primary noise source) to maximise sensitivity to external sound events. The DHT11 sensor was positioned away from the laptop's exhaust fan to avoid heat contamination of temperature readings. The photoresistor was left unshielded to capture ambient room light conditions. A single USB cable connected the board to the laptop for both power and serial data transfer. As all sensors are at a single bedside location, the readings represent conditions at that position rather than across the entire room. The wiring diagram and a photograph of the assembled sensor setup are provided in Appendix~A.

\textbf{Sensor limitations and upgrade path.} Table~\ref{tab:upgrades} summarises the key limitation of each sensor and a recommended upgrade that would address it.

\begin{table}[H]
\centering
\caption{Sensor Limitations and Upgrade Path}
\label{tab:upgrades}
\scriptsize
\begin{tabularx}{\textwidth}{l X X}
\toprule
\textbf{Current Sensor} & \textbf{Limitation} & \textbf{Recommended Upgrade} \\
\midrule
DHT11 & Low accuracy ($\pm2$\,\textdegree C, $\pm5$\%\,RH) and slow response time, limiting the ability to detect rapid changes in bedroom conditions & BME280 ($\pm1$\,\textdegree C, $\pm3$\%\,RH), which also adds barometric pressure sensing \\
\midrule
ELB0604 photoresistor & Outputs only light on/off after manual calibration. Cannot measure how bright the room actually is & BH1750, which outputs light intensity directly as a numeric value in lux (1--65,535\,lx), removing the need for manual calibration \\
\midrule
KY-038 microphone & Outputs a raw voltage with no calibration to decibels or sound frequency, making readings difficult to compare with published noise guidelines & INMP441 I2S digital microphone (24-bit, 61\,dB signal-to-noise ratio, 60\,Hz--15\,kHz range), which would allow sound levels to be measured in decibels and compared against WHO night noise thresholds \\
\bottomrule
\end{tabularx}
\end{table}

% ──────────────────────────────────────────────────────────────────────
\subsection{Data Collection, Storage and Communication}

This section describes how each data source was collected, stored, and communicated to the analysis pipeline. The end-to-end data flow, from raw data sources through to the three CSV files used in Part~2, is illustrated in Appendix~B. The accompanying code for all collection, retrieval, and preprocessing steps can be found in Part~1 of the \texttt{sleep\_analysis} notebook.

\subsubsection{Data Collection and Communication}

\textbf{Bedroom sensors.} The most complex collection path was the bedroom sensors. USB serial was chosen over the Heltec board's built-in WiFi because the laptop was physically close to the sensor board, making a wired connection simpler and more reliable. WiFi would have introduced additional complexity (network configuration, NTP time synchronisation, reconnection logic) without a clear benefit in this setup. No readings were lost to network dropouts over the 14-night study. The data moves through four stages:

\begin{enumerate}[nosep]
    \item Each minute, the Arduino sketch outputs a single CSV-formatted line containing five sensor values over USB serial.
    \item A Python script running in a Jupyter notebook on the laptop captures this serial output, prepends a timestamp from the laptop clock, and writes the result to a CSV file. Each reading is written to disk immediately, so no data is lost if the laptop or notebook crashes mid-night. Timestamps are generated on the laptop rather than the microcontroller, avoiding the complexity of WiFi/NTP synchronisation, which was unnecessary given that the Ultrahuman Ring provides daily (not timestamped) sleep metrics.
    \item Each row in the CSV contains six fields: \texttt{timestamp}, \texttt{temperature\_c}, \texttt{humidity\_pct}, \texttt{light\_\allowbreak detected}, \texttt{sound\_avg}, and \texttt{sound\_peak}. Each night's data is saved as a date-stamped CSV file (e.g.\ \texttt{2026\allowbreak -02-09\_sleep\_\allowbreak data.csv}) in the data folder.
    \item Over the 14-night study, this produced 14 files of approximately 600 rows each. After the collection period, all files were manually combined in Excel into a single \texttt{bedroom\_sensors.csv} for analysis.
\end{enumerate}

\textbf{Air quality.} The air quality data followed a different approach. Rather than fetching data in real-time during each night, which risked missed API calls during overnight sessions, the data was retrieved in bulk from the Breathe London API after the two-week collection period completed. The API was chosen over web scraping as it provides structured JSON that is easier to filter and transform programmatically. The API was found to timeout (HTTP 504) when requesting large date ranges in a single call, so the data was fetched in two weekly chunks with a 15-second pause between requests. The API returns data in long format (one row per measurement per species), which was transformed to wide format (one row per hour) with columns for the hour, PM2.5, and NO\textsubscript{2}. Only readings within the sleep window were retained, and the result was saved as \texttt{breathe\_london\_air\_quality.csv}.

\textbf{Sleep data.} The sleep data was the simplest to collect. The Ultrahuman data was downloaded manually from the website at the end of the collection period and saved as \texttt{ultrahuman\_sleep\_data.csv}.

\subsubsection{Storage and Trade-offs}

All data is stored locally as CSV files on the laptop, with no cloud database or remote server involved. CSV was chosen because the files are human-readable, easy to inspect, and require no database setup. The total data volume is small (8,400 sensor readings, 140 air quality readings, and 14 sleep summaries), making CSV more than adequate for this study. Storing everything locally also keeps the pipeline simple, avoids dependencies on internet connectivity during overnight collection, and ensures personal sleep data remains private.

The trade-off is that there is no automated backup or remote access to the data. If the laptop failed, the collected data could be lost. For a longer study or one involving multiple participants, a cloud-backed database would be worth the added complexity. Trade-off analysis for other factors, including processing, power, and sensor selection, is provided in Appendix~F (Table~\ref{tab:tradeoffs-hw}).

% ──────────────────────────────────────────────────────────────────────
\subsection{Basic Characteristics of the End-to-End Data Acquisition Setup}

Before proceeding to the analysis in Part~2, it is important to understand how the data acquisition system performed in practice, as the sampling resolution, reliability, and latency of each data path affect the confidence that can be placed in the results.

\subsubsection{Sampling Characteristics}

The sampling rate for each source reflects a trade-off between capturing environmental variation and the practical limits of each sensor or API. Table~\ref{tab:sampling} summarises the sampling characteristics and data volume for each source.

\begin{table}[H]
\centering
\caption{Sampling Rates and Data Volume}
\label{tab:sampling}
\scriptsize
\begin{tabularx}{\textwidth}{p{2.4cm} p{1.6cm} >{\raggedleft\arraybackslash}p{1.5cm} >{\raggedleft\arraybackslash}p{1.5cm} X}
\toprule
\textbf{Source} & \textbf{Sampling\newline Rate} & \textbf{Readings/ Night} & \textbf{Total (14 nights)} & \textbf{Justification} \\
\midrule
Bedroom sensors & 1/min & 600 & 8,400 & \textbullet{} Temperature and humidity change slowly, so sub-minute resolution is unnecessary \newline \textbullet{} Sound is sampled at 5\,ms to capture transient noise events (traffic, sirens, door slams) that would be missed at 1-min polling \newline \textbullet{} The resulting minute-level summaries are then logged to the CSV \\
\midrule
Air quality & Hourly & 10 & 140 & Maximum resolution available from the Breathe London API \\
\midrule
Sleep metrics & Daily & 1 & 14 & Output granularity of the Ultrahuman Ring \\
\bottomrule
\end{tabularx}
\end{table}

\subsubsection{Reliability and Data Quality}

Several measures were taken to ensure reliable data capture over the 14-night study. To avoid macOS power management putting the laptop to sleep during overnight collection, the Amphetamine application~\citep{amphetamine2026} was used to keep the system awake while the notebook kernel was active. Minor timing drift between readings is possible due to laptop-clock timestamping, but is limited to a few seconds at most and is negligible relative to the selected sampling rate.

Table~\ref{tab:eval-checks} summarises the evaluation results. All checks passed, confirming no gaps or missing values across the three datasets, so no imputation, interpolation, or exclusion of nights is required before analysis.

\begin{table}[H]
\centering
\caption{Data Collection Evaluation Results}
\label{tab:eval-checks}
\scriptsize
\begin{tabularx}{\textwidth}{p{3.6cm} p{5.0cm} X}
\toprule
\textbf{Check} & \textbf{Source(s)} & \textbf{Result} \\
\midrule
Completeness & \texttt{bedroom\_sensors.csv} & 600 readings per night across all 14 nights (10\,hrs $\times$ 1/min), 100\% complete \\
\midrule
Completeness & \texttt{breathe\_london\_air\_quality.csv} & 10 readings per night across all 14 nights (10\,hrs $\times$ 1/hr), 100\% complete \\
\midrule
Completeness & \texttt{ultrahuman\_sleep\_data.csv} & All 14 rows present, covering the full study period (9--22 Feb 2026) \\
\midrule
Missing values & All three sources & Zero missing values across all columns in all three datasets \\
\midrule
Sleep window validation & \texttt{bedroom\_sensors.csv}, \texttt{breathe\_london\_air\_quality.csv} & All readings fall within the 11\,pm--9\,am sleep window \\
\bottomrule
\end{tabularx}
\end{table}

\subsubsection{Latency}

The analysis is entirely retrospective, so all three sources were merged into a single dataset after the collection period ended. The three data paths therefore operate at different latencies without affecting results: bedroom sensors are captured near real-time via USB serial, air quality data was retrieved in bulk after the study period, and sleep data was downloaded manually at the end of collection. If the system were extended to provide real-time alerts (as proposed in Section~2.4 below), latency requirements would need to be revisited to ensure timely interventions.

% ══════════════════════════════════════════════════════════════════════
\section*{Part 2: Internet of Things}
\setcounter{section}{2}
\setcounter{subsection}{0}
% ══════════════════════════════════════════════════════════════════════

\subsection{Data Interaction/Visualisation/Actuation Platform}

With the data collected, the focus now shifts to analysing it to answer the project's objectives. The \texttt{sleep\_analysis} notebook and the visualisation dashboard serve complementary roles in this: the notebook contains the full statistical pipeline (regression diagnostics, $p$-values, residual tests), while the dashboard provides accessibility and interactivity for a user exploring their own sleep data. Both use the same underlying CSVs and aggregation logic: the notebook discovers the findings, the dashboard communicates them.

The dashboard is organised around the four environment variables with statistically significant Spearman correlations to sleep score (noise variability, humidity, noise level, and temperature range), though all measured variables including light, PM2.5, and NO\textsubscript{2} are still presented for completeness. \textit{As a reminder, all statistical analysis in this study is based on 14 nights of data and should be treated as indicative rather than conclusive.}

\subsubsection{Platform and Technical Implementation}

The dashboard, named RestNet, was built using Streamlit, a Python framework for building web applications, and Plotly, a library for creating interactive graphs. The dashboard and the data files powering it are hosted in a public GitHub repository and deployed automatically via Streamlit Community Cloud. The dashboard reads the same CSV files produced in Part~1 directly from the repository, with all data loading, transformation, and rendering handled by a single Python script.

\subsubsection{User Interface Design}

Three design principles guided the dashboard's interface: a progressive page structure, consistent colour coding, and plain-language readability.

\textbf{Three-page structure.} Rather than presenting all data on a single scrolling page, the dashboard splits content across three pages so that users see summary-level information first and can drill into detail when needed.

\begin{tcolorbox}[
  colback=imperialblue!5,
  colframe=imperialblue,
  leftrule=3pt,
  rightrule=0pt,
  toprule=0pt,
  bottomrule=0pt,
  boxsep=4pt,
  left=8pt,
  right=8pt,
  sharp corners
]
Each page targets a distinct user question:
\begin{itemize}[nosep]
    \item Page~1 (Sleep Dashboard): ``how did I sleep this \textbf{fortnight}?''
    \item Page~2 (Night Explorer): ``what happened on a \textbf{specific night}?''
    \item Page~3 (My Sleep Insights): ``what should I \textbf{change}?''
\end{itemize}
\end{tcolorbox}

\textbf{Colour scheme.} The dashboard uses a dark theme because the primary user context is reviewing sleep data, often in the evening or at night. A dark interface reduces eye strain and feels contextually appropriate, with a warm coral accent providing contrast without harsh brightness. To achieve this level of visual control, all metric cards, section headers, and insight panels were built as custom HTML and CSS rather than using Streamlit's native components, which offered limited control over colour and layout. Two colour systems operate across the dashboard:

\begin{itemize}[nosep]
    \item \textbf{Section colours:} each data source has its own colour (gold for sleep metrics, turquoise for bedroom environment, rose for air quality, purple for comparisons and recommendations), consistent across all three pages so the user always knows what type of content they are viewing.
    \item \textbf{Sleep-score colours:} a separate scheme (green for scores $\geq$80, gold for 70--79, red for $<$70) is applied to bars, dots, and cards across all charts, allowing the user to visually scan any chart and immediately see whether good or poor sleep nights cluster at certain values. These quality tiers are defined by this dashboard and are not official Ultrahuman classifications.
\end{itemize}

\textbf{Readability.} All graphs are interactive: users can hover over any data point to see its exact value, zoom into specific time ranges, and pan across the timeline. All labels use plain language (e.g.\ ``Noise Variability'' rather than ``Std Dev''), targeting a non-technical user who wants to understand their own sleep. Matplotlib was retained in the notebook where static figures are more appropriate.

Further platform and visualisation trade-offs are summarised in Appendix~G (Table~\ref{tab:tradeoffs-vis}).

\subsubsection{Dashboard Pages}

Table~\ref{tab:dashboard-pages} provides a detailed overview of the sections and content presented on each of the three dashboard pages.


\begin{table}[H]
\centering
\caption{Dashboard Page Structure}
\label{tab:dashboard-pages}
\scriptsize
\begin{tabularx}{\textwidth}{p{2.2cm} p{3.2cm} X}
\toprule
\textbf{Page} & \textbf{Section} & \textbf{Content} \\
\midrule
\multirow{3}{*}{\parbox{2.2cm}{1: Sleep\\Dashboard}} & Sleep Overview & \textbullet{} Metric cards (avg score, duration, deep sleep), \textbullet{} Nightly bar chart colour-coded by quality tier \\
\cmidrule{2-3}
 & Bedroom Environment & \textbullet{} Metric cards for the 4 significant predictors, \textbullet{} Tabbed bar charts, each colour-coded by that night's sleep score \\
\cmidrule{2-3}
 & External Air Quality & \textbullet{} PM2.5 and NO\textsubscript{2} line charts, \textbullet{} WHO guideline reference lines (15 and 25\,\textmu g/m\textsuperscript{3}) \\
\midrule
\multirow{5}{*}{\parbox{2.2cm}{2: Night\\Explorer}} & Sleep Architecture & \textbullet{} Sleep score and 2-week average comparison, \textbullet{} Stacked sleep stages bar \\
\cmidrule{2-3}
 & Environment Timeline & \textbullet{} Minute-by-minute line charts (temperature, humidity, noise), \textbullet{} Binary light heatmap \\
\cmidrule{2-3}
 & Night at a Glance & \textbullet{} Bedroom extremes (coolest/warmest, quietest/noisiest hour), \textbullet{} Averages compared to 14-night baseline \\
\cmidrule{2-3}
 & Air Quality Verdict & \textbullet{} Pass/fail cards for PM2.5 and NO\textsubscript{2} against WHO guidelines \\
\cmidrule{2-3}
 & What Was Different? & \textbullet{} 4 significant predictors vs personal medians, \textbullet{} Traffic-light colour coding \\
\midrule
\multirow{4}{*}{\parbox{2.2cm}{3: My Sleep\\Insights}} & Best vs Worst Night & \textbullet{} Side-by-side snapshot cards, \textbullet{} Scores, stage breakdowns, and bedroom conditions \\
\cmidrule{2-3}
 & Correlation Explorer & \textbullet{} Spearman scatter plots for 3 environment predictors and awake time, \textbullet{} Trend lines and $r$ values \\
\cmidrule{2-3}
 & Your Optimal Ranges & \textbullet{} Median-split table, \textbullet{} Good vs poor nights across all 4 significant predictors \\
\cmidrule{2-3}
 & Actionable Takeaway & \textbullet{} Strongest predictor explained with WHO evidence, \textbullet{} Practical recommendations \\
\bottomrule
\end{tabularx}
\end{table}

% ──────────────────────────────────────────────────────────────────────
\subsection{Data Analytics, Inferences and Insights}

This section summarises the statistical analysis used to address the three objectives set out in Section~1.1. The full methodology, including all code, intermediate outputs, and diagnostic checks, is in Part~2 of the \texttt{sleep\_analysis} notebook. The analysis is outlined below in five steps, followed by key findings mapped to the study objectives.

\subsubsection{Analysis Pipeline}

\textbf{Step 1: Nightly feature aggregation.} The three data sources from Part~1 (bedroom sensors, air quality, and sleep metrics) operate at different timescales (minute-level, hourly, and daily respectively), so they must first be aligned to a common unit: one row per night. Since the sleep window spans midnight, each reading is assigned to a ``night'' using a simple rule: readings from midday onward belong to that date's night, while readings before midday belong to the previous date's night. For example, both a 11:30\,pm reading on 9~February and a 1:30\,am reading on 10~February are assigned to the night of 9~February, since they are part of the same sleep session. The same logic is applied to air quality data and serves as the join key across all three sources.

The raw bedroom sensor and air quality data is then converted into 12 nightly environment features:

\begin{itemize}[nosep]
    \item Mean, standard deviation, and range for temperature, humidity, and noise (9 features from bedroom sensors)
    \item Light exposure as a percentage of readings with light detected (1 feature from bedroom sensors)
    \item Mean PM2.5 and NO\textsubscript{2} (2 features from air quality data)
\end{itemize}

\noindent These are merged with the six Ultrahuman sleep metrics into a unified dataset of 14 rows $\times$ 18 columns.

\textbf{Step 2: Spearman correlation analysis.} To identify which environment features are most strongly associated with sleep quality, Spearman rank correlations are computed between all 12 environment features and all 6 sleep metrics. Spearman was chosen over Pearson because normality cannot be verified with $n = 14$~\citep{field2018}. Four features showed statistically significant correlations with sleep score ($p < 0.05$): noise variability, mean humidity, mean noise level, and temperature range. Both noise features were also significantly correlated with total sleep duration, awake time, and light sleep duration, making noise the most consistent predictor across sleep metrics. Light exposure, PM2.5, and NO\textsubscript{2} showed no significant correlations with any sleep metric.

\textbf{Step 3: OLS regression.} To explore whether multiple features together better explain sleep outcomes, OLS regression models were fitted. For each of three target metrics, the two environment features with the highest absolute Spearman correlation were selected as predictors (limited to two by the rule of thumb requiring 7--10 observations per predictor~\citep{field2018}), meaning interactions between more than two features cannot be modelled. All predictor pairs passed the $|r| > 0.7$ multicollinearity threshold~\citep{dormann2013} (Std Noise and Mean Humidity: $r = 0.60$, Std Noise and Mean Noise: $r = 0.57$). Table~\ref{tab:regression} summarises the regression results for all three models.

\begin{table}[H]
\centering
\caption{OLS Regression Results Summary}
\label{tab:regression}
\scriptsize
\begin{tabularx}{\textwidth}{p{3.8cm} X X X}
\toprule
 & \textbf{Sleep Score} & \textbf{Awake Time} & \textbf{Light Sleep} \\
\midrule
Predictors & Noise Variability\newline + Mean Humidity & Noise Variability\newline + Mean Noise & Noise Variability\newline + Mean Humidity \\
\midrule
$F$-statistic ($p$-value) & 12.70 ($p < 0.01$) & 5.84 ($p = 0.019$) & 10.69 ($p = 0.003$) \\
\midrule
Adjusted $R^2$ & 0.643 & 0.427 & 0.598 \\
\midrule
Noise Variability coefficient ($p$-value) & $-4.77$ (0.021) & $+9.17$ (0.050) & $-30.37$ (0.013) \\
\midrule
Second predictor coefficient ($p$-value) & $-0.45$ (0.244) & $+0.63$ (0.283) & $-1.10$ (0.610) \\
\midrule
Durbin--Watson statistic & 2.07 & 1.18 & 1.52 \\
\midrule
Jarque--Bera $p$-value & 0.59 & 0.58 & 0.40 \\
\bottomrule
\end{tabularx}
\end{table}

All three models are statistically significant ($p < 0.05$), with Sleep Score showing the strongest fit (Adj~$R^2 = 0.643$). Noise variability is significant or borderline significant in every model, while the second predictor does not reach significance once noise variability is accounted for. Residual diagnostics (Durbin--Watson, Jarque--Bera, Breusch--Pagan, Shapiro--Wilk) are reported but have limited power with only 14 observations.

\textbf{Step 4: Optimal conditions comparison (median split).} To compare how bedroom conditions differ between good and poor sleep nights, the 14 nights are split at the median sleep score into the top 7 (``good sleep'', scores 86--96) and bottom 7 (``poor sleep'', scores 66--82). A descriptive comparison was used given the low test power at $n = 7$ per group. Mean noise was 15.2 points higher on poor nights (86.0 vs 70.7), noise variability was greater (22.2 vs 20.3), and humidity was 5.3 percentage points higher (53.4\% vs 48.1\%). Temperature differed by only 0.6\,\textdegree C. Light, PM2.5, and NO\textsubscript{2} showed negligible differences (all $<$1 unit).

\textbf{Step 5: Intra-night profiles.} Steps 2--4 use nightly aggregates and cannot capture variation within a single night. To examine how conditions evolved overnight, minute-by-minute sensor data for the best and worst sleep nights was plotted side by side. The worst night was warmer (20.5\,\textdegree C vs 18.6\,\textdegree C), more humid ($\sim$58\% vs $\sim$44\%), and had more frequent noise spikes, particularly around 11\,pm--1\,am and after 5\,am. The best night had early spikes but settled into a quieter, more stable pattern overnight. This illustrates why noise variability (not just average level) emerged as the dominant predictor.

\subsubsection{Key Findings by Study Objective}

\textbf{Objective~1:} to identify which bedroom conditions are most strongly associated with sleep quality. Spearman correlation (Step~2) identified four statistically significant predictors of sleep score: noise variability ($r = -0.80$), humidity ($r = -0.70$), noise level ($r = -0.69$), and temperature range ($r = -0.56$). Light, PM2.5, and NO\textsubscript{2} showed no significant associations. NO\textsubscript{2} and PM2.5 stayed well within WHO guidelines (15\,\textmu g/m\textsuperscript{3} for PM2.5, 25\,\textmu g/m\textsuperscript{3} for NO\textsubscript{2}) throughout the study, so the inability to detect an effect may reflect genuinely clean overnight air during this winter period rather than an absence of any relationship. WHO guidelines are based on 24-hour averages, while this study uses overnight-only data, so comparisons are indicative rather than direct.

\textbf{Objective~2:} to explore whether these associations hold when multiple conditions are considered together. OLS regression (Step~3) showed that when multiple predictors are combined, noise variability was the only factor that retained statistical significance. Humidity and noise level did not reach significance once noise variability was accounted for, suggesting they may be partially confounded with noise patterns rather than acting as independent sleep disruptors. With $n = 14$, the study lacks power to disentangle more subtle interactions.

\textbf{Objective~3:} to compare how bedroom conditions differ between good and poor sleep nights. The median-split comparison (Step~4) showed that the seven poorest sleep nights had substantially higher mean noise, greater noise variability, and higher humidity than the seven best. Intra-night profiling (Step~5) confirmed that the pattern of repeated noise spikes, not sustained loudness, most distinguishes good from poor sleep.

\begin{tcolorbox}[
  colback=imperialblue!5,
  colframe=imperialblue,
  leftrule=3pt,
  rightrule=0pt,
  toprule=0pt,
  bottomrule=0pt,
  boxsep=4pt,
  left=8pt,
  right=8pt,
  sharp corners
]
Overall, \textbf{noise variability} emerged as the factor with the \textbf{strongest and most consistent association} with sleep quality across correlation, regression, group comparison, and intra-night profiling.
\end{tcolorbox}

% ──────────────────────────────────────────────────────────────────────
\subsection{Discussions on the Important Aspects of the Project}

This section reflects on the limitations of the study and considerations around privacy, ethics, and security.

\subsubsection{Limitations and Data Quality}

Several limitations should be considered when interpreting the findings.

\textbf{Sample size and generalisability.} The study involves a single participant over 14 nights, so findings are personalised and not generalisable. With $n = 14$, a Spearman correlation must exceed approximately $|r| = 0.54$ to reach significance at $p < 0.05$, meaning moderate associations (e.g.\ temperature mean, $r = -0.38$) may exist but go undetected. The results should therefore be interpreted as suggestive, offering a sense of which environmental factors may matter most rather than definitive conclusions suitable for clinical or research use.

\textbf{Seasonal and observational constraints.} Data was collected during February in central London, so conditions such as higher temperatures, longer daylight hours, and increased outdoor noise typical of summer months are not represented. For example, summer nights with open windows would likely introduce stronger air quality and noise effects that winter data cannot capture. The observational design also permits no causal claims: humidity and noise level correlated with each other ($r = 0.60$), so it is unclear whether both independently affect sleep or whether one is a proxy for the other. Other unmeasured confounders such as stress, caffeine intake, or screen time before bed could also explain some of the observed associations.

\textbf{Sensor limitations.} The sound sensor has no decibel calibration, so the finding that ``noise variability matters'' cannot be translated into a concrete threshold (e.g.\ ``keep noise below 30\,dB''). Similarly, the binary light sensor cannot distinguish between dim and full room lighting. The Breathe London air quality sensor is located outdoors approximately 500\,m from the bedroom, and indoor pollutant concentrations can differ substantially depending on ventilation and cooking activity. The fixed sleep window (11\,pm to 9\,am) may include pre-sleep activity on nights the participant went to bed late, or miss early wake-ups.

\textbf{Sleep metric dependency.} The Ultrahuman sleep score is a proprietary composite metric whose exact algorithm is not published. This means the study cannot verify how the score weights different sleep stages, and the same raw sleep data might produce different scores on a different wearable platform. Findings are therefore dependent on Ultrahuman's internal calibration and may not correspond directly to clinically validated measures such as polysomnography.

\textbf{Analysis granularity.} The sleep score is a single daily value, so the main analyses (Steps~2--4) use nightly aggregates. A natural next step would be more detailed lag analysis (e.g.\ ``does a noise spike at 2\,am predict more awake time in the following hour?''), which would require minute-level sleep stage data that the Ultrahuman export does not currently provide.

\subsubsection{Privacy, Ethics, and Security}

\textbf{Privacy.} The study collects two categories of personal data: sleep metrics from a wearable device (sleep score, duration, stage breakdown) and continuous bedroom environment readings (temperature, humidity, noise, light). Under the UK GDPR~\citep{ukgdpr2018}, sleep data qualifies as health-related personal data (Article~9), which requires both a lawful basis for processing and a specific condition for handling special category data. As the sole participant and data subject, the researcher provided implicit consent for processing their own data, satisfying the explicit consent condition under the regulation. A multi-user deployment would require written informed consent specifying what data is collected, how long it is retained, and the right to withdraw. Several steps were taken to minimise privacy risk: all data is stored locally as CSV files with no cloud database or third-party storage, the dashboard displays the participant as ``User Y.A.'' rather than by full name, and the public GitHub repository uses pseudonymised data with no direct identifiers beyond dates and the repository username. Full anonymisation would require removing dates and hosting under a non-identifying account.

\textbf{Ethics.} Continuous bedroom monitoring raises ethical considerations even in a self-study. The sound sensor records noise levels (analog voltage) rather than audio, avoiding the more intrusive concerns of bedroom audio capture. However, if deployed in shared households, the system would passively monitor the sleep environment of other occupants who have not consented. The ICO's draft guidance on consumer IoT products~\citep{ico2025iot} notes that IoT devices can process special category data including health information, and that data protection law applies to the manufacturer even when users process data in a domestic setting. A multi-participant study would also require a clear data retention policy (e.g.\ deletion after analysis) and transparency about how findings are used.

\textbf{Security.} Data is currently stored as unencrypted CSV files on a single laptop, which presents a risk if the device is lost or stolen. The Streamlit dashboard is served over HTTPS (provided by Streamlit Community Cloud), but there is no user authentication, meaning anyone with the URL can view the data. For a production deployment, the following measures would be needed: encrypted storage at rest (e.g.\ AES-256), user authentication on the dashboard, secure API endpoints if sensor data were transmitted wirelessly, and regular backups to prevent data loss. The current USB serial connection avoids wireless interception risks.

% ──────────────────────────────────────────────────────────────────────
\subsection{Avenues for Future Work and Potential Impact}

The findings and limitations discussed in Section~2.3 point to several directions for improving and extending this study, ranging from straightforward hardware upgrades to broader applications.

\subsubsection{Future Extensions}

\textbf{Sensor upgrades.} Replacing the current temperature, humidity, light, and sound sensors with calibrated alternatives (e.g.\ decibel-rated microphone, lux-capable light sensor) would allow direct comparison with published health guidelines such as the WHO night noise guideline of $<$40\,dB L\textsubscript{night,outside}~\citep{who2009}. Beyond improving existing measurements, the kit could be extended with entirely new sensors: a CO\textsubscript{2} sensor (e.g.\ SCD40) would add a ventilation proxy, since elevated CO\textsubscript{2} in closed bedrooms has been associated with reduced sleep quality~\citep{basner2023, caddick2018}, and a barometric pressure sensor could potentially capture weather-related changes that may affect sleep.

\textbf{Longer study and multiple participants.} Extending the study to three months ($n \approx 90$ nights) would help capture subtler effects such as temperature or light, which may require larger samples to reach statistical significance. Depending on timing, it could also cover seasonal variation (e.g.\ winter into spring) and support more predictors per regression model without overfitting. Including multiple participants would help test whether noise variability is a broadly relevant sleep disruptor or specific to this individual and location, and could enable between-subject comparisons across different bedroom setups and neighbourhoods.

\textbf{Real-time alerts and smart home integration.} The current system is open-loop: it records conditions but cannot intervene. The Heltec board's built-in WiFi or LoRa radio could close this loop by pushing alerts when conditions exceed the good-sleep group means identified in Step~4 (e.g.\ noise variability rising above 20.3, humidity exceeding 48\%). This could extend to smart home actuation: triggering a white noise machine when noise spikes are detected, or a dehumidifier when humidity rises, creating a closed-loop system (sense $\rightarrow$ assess $\rightarrow$ intervene $\rightarrow$ measure outcome).

\textbf{Machine learning prediction.} With sufficient data (e.g.\ 90+ nights from a longer study), a sleep quality classifier such as random forest or gradient boosting could be trained on the 12 environment features to predict tonight's expected sleep quality from current bedroom conditions. This would move the system from retrospective analysis to prospective guidance, alerting the user before bed that conditions are suboptimal. Scaling from one bedroom to many introduces practical challenges, however: different sensor models may need recalibration, sleep trackers from different brands produce scores on different scales, and a model trained on one person's data may not generalise to another. Handling these differences reliably becomes a significant engineering problem as the system grows.

\subsubsection{Broader Impact}

The total hardware cost of this system is under \pounds30, compared with clinical polysomnography studies that typically cost over \pounds1{,}000 per night per participant. While low-cost sleep environment sensors are not new, this combination of price and open-source design lowers the barrier to repeated, multi-night monitoring at individual or household level. The full codebase is public, so anyone could replicate the setup without specialist equipment. This also opens practical commercial routes: a landlord could offer tenants environment dashboards as part of a building wellness programme, or a sleep coaching service could use the sensor kit as a starter pack for clients.

The full hardware design, Arduino sketch, analysis notebook, and dashboard code are published as an open-source repository, enabling replication of the sensing and analysis pipeline, though a wearable sleep tracker (or substitute) would still be needed for the sleep quality outcome data. A neighbourhood-scale deployment, for example 20 households each running the same sensor kit for three months, would yield sufficient data to detect smaller effects and control for between-household variation in building insulation, floor level, and street-facing orientation. Such a dataset could test whether the noise variability effect observed here holds across different urban environments or is specific to this participant's location and sleeping habits.

% ── Conclusion ────────────────────────────────────────────────────────
\section*{Conclusion}

This study set out to investigate which bedroom environment conditions are most strongly associated with sleep quality. Spearman correlations identified four significant predictors: noise variability, humidity, noise level, and temperature range. OLS regression showed that noise variability was the only predictor that retained significance when combined with others. Median-split and intra-night comparisons confirmed that poor sleep nights were consistently noisier, more humid, and characterised by more frequent noise spikes.

Across all four analyses, noise variability emerged as the factor with the strongest and most consistent association with sleep quality. Reducing noise variability, through earplugs, white noise machines, or improved window insulation, is the single most impactful change indicated by this dataset.

The end-to-end system worked well for a 14-night observational study. The sensor hardware ran reliably every night with no data loss, the analysis pipeline produced reproducible results, and the dashboard is designed to communicate the findings to a non-technical audience. The main limitation was not the system itself but the study duration: 14 nights provides enough data to identify strong patterns but not enough to confirm them statistically. In hindsight, choosing a calibrated sound sensor and a lux-capable light sensor from the start would have made the findings more directly comparable with published guidelines.

\begin{tcolorbox}[colback=imperialblue!4, colframe=imperialblue, coltext=black, boxrule=0pt, leftrule=3pt, arc=0pt, left=8pt, right=8pt, top=6pt, bottom=6pt]
Overall, this study demonstrates that \textbf{affordable IoT hardware} (under \pounds30), combined with \textbf{public APIs} and a \textbf{consumer wearable}, can produce \textbf{actionable, personalised sleep environment insights}.
\end{tcolorbox}

% ── References ────────────────────────────────────────────────────────
\newpage
\begin{thebibliography}{17}

\bibitem[Amphetamine(2026)]{amphetamine2026}
Amphetamine (2026) \textit{Amphetamine for macOS} [Online]. Available at: \url{https://apps.apple.com/us/app/amphetamine/id937984704} (Accessed: March 2026).

\bibitem[Basner et~al.(2023)]{basner2023}
Basner, M., Smith, M.G., Jones, C.W. et al. (2023) `Associations of bedroom PM2.5, CO2, temperature, humidity, and noise with sleep: an observational actigraphy study', \textit{Sleep Health}, 9(3), pp.~253--263. doi:10.1016/j.sleh.2023.02.010.

\bibitem[Breathe~London(2026)]{breathelondon}
Breathe London (2026) Air Quality Monitoring Network. Available at: \url{https://www.breathelondon.org/} (Accessed: March 2026).

\bibitem[Caddick et~al.(2018)]{caddick2018}
Caddick, Z.A., Gregory, K., Arsintescu, L. and Flynn-Evans, E.E. (2018) `A review of the environmental parameters necessary for an optimal sleep environment', \textit{Building and Environment}, 132, pp.~11--20. doi:10.1016/j.buildenv.2018.01.020.

\bibitem[Dormann et~al.(2013)]{dormann2013}
Dormann, C.F. et al. (2013) `Collinearity: a review of methods to deal with it and a simulation study evaluating their performance', \textit{Ecography}, 36(1), pp.~27--46. doi:10.1111/j.1600-0587.2012.07348.x.

\bibitem[Field(2018)]{field2018}
Field, A. (2018) \textit{Discovering Statistics Using IBM SPSS Statistics}. 5th edn. London: SAGE Publications.

\bibitem[Halperin(2014)]{halperin2014}
Halperin, D. (2014) `Environmental noise and sleep disturbances: A threat to health?', \textit{Sleep Science}, 7(4), pp.~209--212. doi:10.1016/j.slsci.2014.11.003.

\bibitem[ICO(2025)]{ico2025iot}
Information Commissioner's Office (2025) \textit{Guidance for consumer Internet of Things products and services}. Available at: \url{https://ico.org.uk/for-organisations/uk-gdpr-guidance-and-resources/online-tracking/guidance-for-consumer-internet-of-things-products-and-services/} (Accessed: March 2026).

\bibitem[Medic et~al.(2017)]{medic2017}
Medic, G., Wille, M. and Hemels, M.E. (2017) `Short- and long-term health consequences of sleep disruption', \textit{Nature and Science of Sleep}, 9, pp.~151--161. doi:10.2147/NSS.S134864.

\bibitem[Muzet(2007)]{muzet2007}
Muzet, A. (2007) `Environmental noise, sleep and health', \textit{Sleep Medicine Reviews}, 11(2), pp.~135--142. doi:10.1016/j.smrv.2006.09.001.

\bibitem[ResMed(2025)]{resmed2025}
ResMed (2025) \textit{ResMed's fifth annual global sleep survey reveals a world struggling with poor sleep}. Available at: \url{https://newsroom.resmed.com/news-releases/news-details/2025/ResMeds-Fifth-Annual-Global-Sleep-Survey-Reveals-a-World-Struggling-with-Poor-Sleep/} (Accessed: March 2026).

\bibitem[Sleep~Cycle(2025)]{sleepcycle2025}
Sleep Cycle (2025) \textit{New Sleep Cycle report: 105 million nights reveal a global decline in sleep quality}. Available at: \url{https://sleepcycle.com/newsroom/press-release/new-sleep-cycle-report-unveiling-global-sleep-trends-and-insights} (Accessed: March 2026).

\bibitem[UK~GDPR(2018)]{ukgdpr2018}
UK GDPR (2018) \textit{The General Data Protection Regulation as incorporated into UK law by the European Union (Withdrawal) Act 2018 and the Data Protection, Privacy and Electronic Communications (Amendments etc) (EU Exit) Regulations 2019}. Available at: \url{https://www.legislation.gov.uk/eur/2016/679/contents} (Accessed: March 2026).

\bibitem[Ultrahuman(2026)]{ultrahuman}
Ultrahuman (2026) Ultrahuman Ring Air. Available at: \url{https://www.ultrahuman.com/ring/} (Accessed: March 2026).

\bibitem[WHO(2009)]{who2009}
World Health Organization (2009) \textit{Night Noise Guidelines for Europe}. Copenhagen: WHO Regional Office for Europe. Available at: \url{https://www.who.int/europe/publications/i/item/9789289041737} (Accessed: March 2026).

\bibitem[Yasmeen et~al.(2025)]{yasmeen2025}
Yasmeen, S., Li, B., Du, C. and Liu, H. (2025) `Exploring the interconnection of sleep quality, indoor environmental factors, and energy efficiency: strategies for sustainable sleep environments', \textit{Indoor Air}, 2025, 8245786. doi:10.1155/ina/8245786.

\bibitem[Zhao et~al.(2024)]{zhao2024}
Zhao, Y., Zhang, S., Guo, L. et al. (2024) `Association between air pollutants and the risk of sleep disorders: a systematic review and meta-analysis', \textit{Aerosol and Air Quality Research}, 24, 230197. doi:10.4209/aaqr.230197.

\end{thebibliography}

% ── Appendix ──────────────────────────────────────────────────────────
\newpage
\section*{Appendix: Figures and Illustrations}

\textit{This appendix contains all figures, diagrams, and screenshots referenced in the report.}

\subsection*{Appendix A: Wiring Diagram and Sensor Assembly}

\begin{figure}[H]
\centering
\begin{minipage}[t]{0.48\textwidth}
\centering
\includegraphics[width=\textwidth]{figures/wiring_diagram.png}
\par\vspace{4pt}
\textit{\small (a) Circuit diagram: Heltec WiFi LoRa 32 V3 connected to KY-015/DHT11 (GPIO~7), ELB0604 photoresistor (GPIO~1), and KY-038 sound sensor (GPIO~2), all at 3.3\,V.}
\end{minipage}%
\hfill
\begin{minipage}[t]{0.48\textwidth}
\centering
\includegraphics[width=\textwidth]{figures/sensor_assembly.png}
\par\vspace{4pt}
\textit{\small (b) Sensor assembly: KY-015/DHT11 (left), KY-038 microphone (middle), ELB0604 photoresistor (right) on a mini breadboard, powered via USB.}
\end{minipage}
\end{figure}

\vspace{12pt}

\noindent
\begin{minipage}[t]{0.48\textwidth}
\subsection*{Appendix B: Data Flow Diagram}
\textit{\small End-to-end data pipeline: three sources (Heltec board, Breathe London API, Ultrahuman app) producing CSV files, unified into a single dataset for notebook analysis and Streamlit dashboard.}

\vspace{4pt}
\centering
\includegraphics[width=\textwidth]{data_flow_pipeline.png}
\end{minipage}%
\hfill
\begin{minipage}[t]{0.48\textwidth}
\subsection*{Appendix C: Spearman Correlation Heatmap}
\textit{\small Heatmap showing Spearman rank correlations between all 12 environment features and 6 sleep metrics, with significance stars ($^*\,p < 0.05$, $^{**}\,p < 0.01$).}

\vspace{4pt}
\centering
\includegraphics[width=\textwidth]{spearman_corr_heatmap.png}
\end{minipage}

\clearpage
\subsection*{Appendix D: Good vs Poor Sleep Nights}
\vspace{-2pt}
\textit{\small Bar charts comparing environment conditions (mean and standard deviation) on the top~7 vs bottom~7 nights by sleep score.}

\begin{center}
\includegraphics[width=0.65\textwidth]{figures/good_vs_poor_bars.png}
\end{center}

\vspace{-4pt}
\subsection*{Appendix E: Intra-Night Environment Profiles}
\vspace{-2pt}
\textit{\small Side-by-side minute-by-minute plots of temperature, humidity, and noise for the best and worst sleep nights.}

\begin{center}
\includegraphics[width=0.65\textwidth]{figures/intra_night_profiles.png}
\end{center}

\subsection*{Appendix F: Hardware and System Trade-offs}

\begin{table}[H]
\centering
\caption{Hardware and System Trade-offs}
\label{tab:tradeoffs-hw}
\scriptsize
\setlength{\tabcolsep}{4pt}
\begin{tabular}{@{} p{2cm} p{2.8cm} >{\centering\arraybackslash}p{0.8cm} p{5.6cm} p{5.2cm} @{}}
\toprule
\textbf{Trade-off} & \textbf{Options (A / B)} & \textbf{Chosen} & \textbf{Gain} & \textbf{Cost} \\
\midrule
Processing & Local / Cloud & A & No internet dependency overnight, zero latency, personal sleep data stays on device & No remote access, no automated backup \\[4pt]
\midrule
Communication & USB serial / WiFi & A & Simpler, no network config, zero dropped readings over 14 nights & Laptop must be physically close, no remote monitoring \\[4pt]
\midrule
Power & Continuous USB / Battery & A & No battery constraint over 10-hr sleep windows, always-on & Not portable, tied to laptop location \\[4pt]
\midrule
Sensors & Low-cost (${\sim}$\pounds30) / Calibrated (${\sim}$\pounds60+) & A & 4$\times$ cheaper, faster prototyping, sufficient to detect relative patterns & No absolute dB or lux; DHT11 $\pm2$\,\textdegree C vs BME280 $\pm1$\,\textdegree C \\[4pt]
\midrule
Analysis & Spearman/OLS / ML & A & Transparent, interpretable coefficients, works with $n{=}14$ & May miss non-linear relationships; justified by small sample \\[4pt]
\midrule
Data privacy & All local / Cloud & A & Zero exposure of personal sleep data to third parties & No collaborative features, no multi-device sync \\
\bottomrule
\end{tabular}
\end{table}

\subsection*{Appendix G: Platform and Visualisation Trade-offs}

\begin{table}[H]
\centering
\caption{Platform and Visualisation Trade-offs}
\label{tab:tradeoffs-vis}
\scriptsize
\setlength{\tabcolsep}{4pt}
\begin{tabular}{@{} p{2cm} p{2.8cm} >{\centering\arraybackslash}p{0.8cm} p{5.6cm} p{5.2cm} @{}}
\toprule
\textbf{Trade-off} & \textbf{Options (A / B)} & \textbf{Chosen} & \textbf{Gain} & \textbf{Cost} \\
\midrule
Functionality vs.\ Rapid development & Streamlit / Flask with Dash & A & Rapid prototyping with built-in widgets, reactive model, and free cloud hosting & Limited layout control; script re-executes on every user interaction, requiring explicit caching and state management \\[4pt]
\midrule
Price vs.\ Robustness & Streamlit Community Cloud (free) / Self-hosted server & A & Zero cost, automatic deploys from GitHub, HTTPS by default & Limited server resources, no custom domain, cold starts after inactivity \\[4pt]
\midrule
Performance vs.\ Usability & Cached data loading / No caching & A & Data processing runs once per session rather than on every page change, keeping navigation responsive & Slightly higher memory usage; cached results not updated if source files change mid-session \\[4pt]
\midrule
Functionality vs.\ Usability & Full statistical output (notebook) / Simplified visual summaries (dashboard) & Both & Notebook provides full regression diagnostics and p-values for academic rigour; dashboard presents key findings in plain language with interactive charts for non-technical users & Two separate codebases to maintain; aggregation logic duplicated across notebook and dashboard script \\[4pt]
\midrule
Price vs.\ Usability & Custom HTML/CSS theming / Streamlit native components & A & Full control over colour scheme, section headers, and metric card layout, enabling sleep-score colour coding and dark theme & 88 \texttt{unsafe\_allow\_html=True} calls; higher development time; harder to maintain if Streamlit updates break HTML rendering \\[4pt]
\midrule
Latency/Delay vs.\ Accuracy & Real-time per-reading analysis / Retrospective batch analysis & B & All three sources merged after collection, ensuring complete data with no gaps before any statistics are computed & No real-time feedback during the study; user cannot act on conditions as they happen overnight \\
\bottomrule
\end{tabular}
\end{table}

\newpage
\subsection*{Appendix H: Sleep Dashboard Page}
\vspace{-4pt}
\textit{\small Sleep score bar chart, bedroom environment tabs, and air quality trends. Interactive version: \url{https://sleep-monitor.streamlit.app}.}
\vspace{2pt}
\begin{center}
\includegraphics[width=\textwidth,height=0.88\textheight,keepaspectratio]{figures/dashboard_sleep.pdf}
\end{center}

\newpage
\subsection*{Appendix I: Night Explorer Page}
\vspace{-4pt}
\textit{\small Sleep architecture, sensor timelines, night-at-a-glance summary, air quality verdict, and predictor comparisons. Interactive version: \url{https://sleep-monitor.streamlit.app}.}
\vspace{2pt}
\begin{center}
\includegraphics[width=\textwidth,height=0.88\textheight,keepaspectratio]{figures/dashboard_explorer.pdf}
\end{center}

\newpage
\subsection*{Appendix J: My Sleep Insights Page}
\vspace{-4pt}
\textit{\small Best vs worst night comparison, correlation scatter plots, optimal ranges, and actionable takeaway. Interactive version: \url{https://sleep-monitor.streamlit.app}.}
\vspace{2pt}
\begin{center}
\includegraphics[width=\textwidth,height=0.88\textheight,keepaspectratio]{figures/dashboard_insights.pdf}
\end{center}

\end{document}
